\documentclass{article}

%!TEX root = ../main.tex

\usepackage{url}
\def\UrlFont{\rmfamily}

\usepackage{enumitem}
% [noitemsep,topsep=0pt,labelwidth=0pt,listparindent=-30pt,leftmargin=30pt,rightmargin=30pt]
\usepackage{cite}
\usepackage{amsmath,amssymb,amsfonts}
\usepackage{algorithmic}
\usepackage{graphicx}
\graphicspath{{./img/}}

\usepackage{textcomp}
\usepackage{xcolor}
\usepackage{booktabs}
\usepackage{caption}
\usepackage{xspace}
\usepackage{float}
\usepackage{marginnote}
\usepackage{tabularx}
\usepackage{changepage, threeparttable}
\usepackage{attrib}
\usepackage{svg}
\usepackage{multirow}


\graphicspath{{./img/}}
%!TEX root = ../main.tex


\newcommand\cc[1]{\textcolor{red}{#1}}
\newcommand\tobedeleted[1]{\textcolor{red}{#1}}



\title{Scrum Simulation %
	\normalsize{[v26.1]}}

\date{\today}
\begin{document}
\maketitle


\vfill
\begin{description}[noitemsep,topsep=0pt,labelwidth=0pt,listparindent=-30pt,leftmargin=30pt,rightmargin=30pt]
\item[Team name: ] Group 16 
\item[Company: ] Alfa Laval  
\item[Teaching assistant: ] Hiep Cao
\item[Agile coach: ] Henning Therkelsen
\end{description}
\begin{table}[h]\begin{center}
\begin{tabular*}{1\textwidth}{@{}ccp{8cm}}
\toprule
Role & email & Name \\
\midrule
PO & Kasper.Haarslov@alfalaval.com & Kasper Haarsløv \\
Specialist & morten.gronkjaer@alfalaval.com & Morten Groenkjaer \\
SM & maev@itu.dk & Maria Melina Evdaimon \\
dev & cmje@itu.dk & Casper Meier Jenkins \\
dev & ebll@itu.dk & Emilie Bliddal Ravn Larsen \\
dev & jpre@itu.dk & Joakim-David Stauning \\
dev & kfje@itu.dk & Kasper Fjord Grønvold Jensen \\
dev & mhil@itu.dk & Milja Noomi Hildigunnsdottir \\
dev & moto@itu.dk & Morgan Møller Torp \\
dev & merl@itu.dk & Morten Holtorp Erlandsen \\
dev & thsch@itu.dk & Therese Schultz-Nielsen \\
dev & vmol@itu.dk & Vincent Jørss Molenaar\\
\bottomrule
\end{tabular*}\end{center}\end{table}


\newpage



\section{After the Scrum simulation}
\subsubsection*{Tasks}
\begin{itemize}
    \item Execute the test at \url{https://www.scrum.org/open-assessments/scrum-open}.
    \item Percentage: 70\% 
    \item "When does the next Sprint begin?"
        Here I answered that it started after the next Sprint Planning but was surprised that it begins immediately after

        "Who is required to attend the Daily Scrum?"
        I Thought that the Scrum Master was supposed to attend the Daily Scrum
\end{itemize}



\subsubsection*{What are the main points I learned through the simulation?}
\begin{itemize}
    \item In short time frame, 30 min per sprint, Scrum is wasteful 
    \item It is hard to meet the demand in 4+4 min 
    \item It was hard to estimate when we did not know the work time 
\end{itemize}

\subsubsection*{Share overall observations about your experience as a participant of the simulation.}

I did not like the LEGO scrum simulation because it was a stressful event. 
We could not use our estimates since they were created with wrong assumptions and we did not have the time to re-estimate.
It was hard to build anything with a 4+4 minutes work with a 2 min meeting.
In the simulation, because of the time limit, every meeting, Daily Scrum, Scrum- Retrospective and Planning, seemed wasteful and stressful.
It also was a little hard to work with a split PO since she needed to be at both Scrum teams.
This with the limited time frame further added to the stressfulness of the event.
A problem with the event was when we slightly ignored Scrum we seemed to get work done.
This I see more as a problem with us only having 4+4 minutes to build and not representative of Scrum.
Because of this I do not think there is much, if something at all, I can take away from the event.

\end{document}
